% -----------------------------------------------------------------------------
% Section 5: Implications
% Paper 5 — Constitutional Code (Separation of Powers)
% Security audit patch applied 2026-03-29:
%   Patch 5.2: Residual semantic gap paragraph after trade-secret dissolution
% -----------------------------------------------------------------------------

\section{Implications}
\label{sec:implications}

\subsection{For AI Governance}

The separation-of-powers architecture has three immediate implications
for regulated AI deployments.

First, it resolves the trade-secret tension without disclosure.
Operators routinely resist algorithmic transparency requirements on
the grounds that model weights, training data, and scoring rules
constitute trade secrets. The Legislative--Executive--Judicial
tripartite structure dissolves this tension: the Judicial branch can
verify that decisions are consistent with the declared policy
(Legislative) and were produced by an unmodified kernel (Executive)
without ever accessing the model internals. The ZK-SNARK layer of
Paper~3~\cite{btv-ar2026} makes this precise: an auditor can verify
$\varepsilon$-statistical consistency of the operator's redaction
history without learning which specific verdicts were redacted or why.
The architecture renders disclosure unnecessary for verification.

\paragraph{Residual semantic gap.}
The ZK verification guarantee above holds under the assumption that
the demographic attributes committed to the transparency log at
ingestion time are truthful. If the operator commits false group
labels, the ZK proof certifies compliance with a fabricated
statistical history---a semantically vacuous guarantee. Paper~3
(Section~7, Limitation~3 of~\cite{btv-ar2026}) identifies this as
the \emph{garbage-in} vulnerability and proposes an Attested
Demographics extension whereby attribute commitments require a
digital signature from a trusted external source (civil registry,
identity provider). Until this extension is implemented, the
Judicial branch verifies \emph{structural} fairness (the operator
cannot selectively erase evidence) but not \emph{substantive}
fairness (the operator's initial data was correct). This mirrors
the semantic-truthfulness boundary of Paper~1
(Section~6.1 of~\cite{btv2026}): the type system guarantees
that evidence exists, not that the evidence is true. Closing this
gap fully requires institutional mechanisms---accredited data
sources, cross-referencing with external registries---that
complement the cryptographic architecture.

Second, it operationalises Montesquieu's separation of powers as a
software architecture. The Legislative branch (policy YAML in version
control) cannot execute decisions; the Executive branch (the Rust
kernel) cannot modify policy; the Judicial branch (ZK auditors) can
verify both without participating in either. This mirrors the
constitutional design principle that accountability requires not merely
transparency but structural impossibility of self-dealing.

Third, it provides a concrete compliance path for the EU AI Act
(Art.~13 transparency, Art.~14 human oversight, Art.~17 quality
management). The Constitutional Completeness Corollary
(Corollary~\ref{cor:completeness}) guarantees that the three branches
are capability-disjoint: no single actor can simultaneously legislate,
execute, and adjudicate---the structural precondition for the
accountability requirements of Art.~9.

\subsection{For Software Engineering}

The constitutional metaphor is more than an analogy: it is a design
pattern for accountability-critical systems. Rust's type system
provides the enforcement mechanism that makes the separation
structurally guaranteed rather than merely conventional. The pattern
generalises beyond AI governance: any system where a single actor's
unilateral action could cause irreversible harm---financial settlement,
medical authorisation, criminal sentencing---benefits from the same
architectural discipline.
